\section{Reproducibility}
\label{sec:reproducibility}

The repository is organized so that every number in this report can be traced
to an audited artifact:

\begin{itemize}
  \item \textbf{Release artifacts.} The seven JSON files under
    \texttt{results/final/} are the authoritative evidence. Each records the
    Git commit (\texttt{0a11fda}, \texttt{dirty=false}), SHA-256 hashes of the
    embedding and qrels inputs, package versions, CPU affinity and thread
    environment, warm-up and repetition counts, raw outer-timer samples,
    per-stage timings, exact-ranking recovery, qrels metrics, and, for BOND,
    per-checkpoint pruning counts and score-validation errors.
    \texttt{results/final/manifest.json} lists the accepted files with sizes
    and hashes. Historical results under \texttt{results/legacy/} and
    dirty-worktree pilots under \texttt{results/controlled/} are preserved for
    auditability but support no final performance claim.
  \item \textbf{Pinned environments.} Python dependency manifests are pinned;
    the PDX library is checked out at the audited commit \texttt{fdc62f2}
    (the controlled runner refuses a dirty or unexpected PDX checkout).
    Key package versions recorded in the artifacts include NumPy 2.4.6,
    faiss-cpu 1.14.2, PyLate 1.5.0, fast-plaid 1.3.0.290, and CPU
    PyTorch 2.9.0 on Python 3.11.9.
  \item \textbf{Independent kernels and tests.} The compiled exact kernel
    (\texttt{cpp/exact\_maxsim/}) and the exact-safe BOND kernel
    (\texttt{cpp/bond\_maxsim/}) ship with build scripts and automated
    unit tests against the NumPy oracle, including tie, all-negative,
    forced-pruning, and invalid-input cases.
  \item \textbf{Controlled runners.} \texttt{experiments/pipeline/08}
    (IVF and PLAID), \texttt{10} (BOND), \texttt{13} (validation-only PLAID
    selection), and \texttt{14} (PLAID figure generation) implement the protocol of
    \cref{sec:protocol}; the repository \texttt{README.md} documents the
    exact \texttt{taskset} and environment-variable invocations used for the
    release runs.
  \item \textbf{Figures from artifacts.} Figures
    \ref{fig:quality-latency}--\ref{fig:bond} are generated directly from the
    release JSON files by \texttt{experiments/pipeline/09}, \texttt{11},
    \texttt{12}, and \texttt{14}, without hard-coded headline values.
\end{itemize}

Large embedding inputs remain outside Git; their SHA-256 hashes inside each
artifact tie the recorded results to immutable inputs. The documentation set
(\texttt{docs/benchmark\_protocol.md}, \texttt{docs/audit\_findings.md},
\texttt{docs/bond\_kernel\_design.md}, \texttt{docs/final\_results.md})
records the protocol, the audit, the kernel safety argument, and the release
summary, respectively.
